%\documentclass[11pt]{scrartcl}
\documentclass[11pt, twoside, a4paper, BCOR8mm, DIV12, bibtotoc,idxtotoc]{scrbook}
\usepackage{german}
\usepackage{typearea}
\usepackage{longtable}
\usepackage{hyperref}
\usepackage{graphicx}

% Zusaetzliche Picture-Umgebungen (z.B. shadowenv)
\usepackage{picins}

% Header anpassbar
\usepackage{fancyhdr}

% Headings umdefinieren
\pagestyle{fancy}
\fancyhf{}
\fancyhead[RO]{\nouppercase{\rightmark}}
\fancyhead[LE]{\nouppercase{\leftmark}}
\fancyfoot[RO, LE]{\thepage}

%\addtolength{\headwidth}{\marginparsep}
%\addtolength{\headwidth}{\marginparwidth}
\addtolength{\headwidth}{1cm}

\parindent0.0mm
\parskip0.3cm    
\typearea{13}

\begin{document}

\frontmatter

\begin{titlepage}

\begin{center}
\rule[-.1in]{16cm}{1mm}\\[3mm]
{\fontfamily{cmss}\fontseries{bx}\fontshape{n}\fontsize{20}{20pt}\selectfont
  www.openbib.org $\bullet$ OpenBib Rechercheportal}\\[-2mm]
\rule[-.1in]{16cm}{1mm}

\vspace{5cm}

  \textbf{\fontfamily{cmss}\fontseries{bx}\fontshape{n}\fontsize{30}{30pt}\selectfont OpenBib Recherche-Portal\\[3mm] Installations-Handbuch}

  \vspace{2cm}

  Oliver Flimm \texttt{<flimm@openbib.de>}\\
  Stand: \today

  \vspace{8cm}

\rule[-.1in]{16cm}{1mm}\\[3mm]
{\fontfamily{cmss}\fontseries{bx}\fontshape{n}\fontsize{20}{20pt}\selectfont
  www.openbib.org $\bullet$ OpenBib Rechercheportal}\\[-2mm]
\rule[-.1in]{16cm}{1mm}

\end{center}

\end{titlepage}

%\thispagestyle{empty}

%\begin{verbatim}


%Copyright (c) 2004 Oliver Flimm <flimm@openbib.org>

%Es wird die Erlaubnis gegeben dieses Dokument zu kopieren, verteilen 
%und/oder zu veraendern unter den Bedingungen der GNU Free
%Documentation License, Version 1.1 oder einer spaeteren, von der Free 
%Software Foundation veroeffentlichten Version; mit den
%Unveraenderlichen Abschnitten DEREN TITEL AUFGEZAEHLT sind, mit den 
%Vorderseitentexten die AUFGEZAEHLT sind, und mit den Rueckseitentexten
%die AUFGEZAEHLT sind. Eine Kopie dieser Lizenz ist in dem Abschnitt 
%enthalten, der mit "GNU Free Documentation License"
%\end{verbatim}

\tableofcontents

\mainmatter

\chapter{Einleitung}

Das OpenBib-Rechercheportal besteht aus verschiedenen Komponenten. Es
sind dies

\begin{itemize}
\item das eigentliche Recherche-Portal als Web-Anwendung
\item die darunterliegenden Katalog-, Session und Benutzer-Datenbanken
  sowie
\item die Konvertierungsprogramme, um die Katalog-Datenbanken 
  automatisiert aus den externen Daten verschiedener
  (Bibliotheks)Systeme zu f"ullen.
\end{itemize}

Alle Komponenten setzen eine Infrastruktur voraus. Konkret besteht
diese aus einem Rechner mit Betriebssystem, Webserver, Datenbanksystem
sowie einer Programmierumgebung. Entwickelt wird OpenBib mit dem
Betriebssystem Linux (Debian 'sarge'), dem Webserver Apache (mit
mod\_perl), dem DBMS MySQL sowie verschiedenen externen Perl-Modulen.


Um das Recherche-Portal zu installieren sind folgende Schritte n"otig:
\begin{enumerate}
\item Download der OpenBib-Software
\item Installation der notwendigen Infrastruktur
\item Installation von OpenBib und seinen Komponenten
\end{enumerate}

Die Installation wird in diesem Handbuch beispielhaft mit folgenden
Festlegungen durchgef"uhrt:

\begin{itemize}
\item Das verwendete Betriebssystem ist Linux und die verwendete
  Distribution ist Debian 3.1 ('sarge').
\item Es wird lediglich ein Server f"ur OpenBib verwendet. Als Name
  wird in der Dokumentation \texttt{openbib.foo.bar} verwendet.
\item Die Perl-Module der OpenBib-Webanwendung werden unter
\begin{verbatim}
/usr/local/lib/site_perl/
\end{verbatim}
  installiert.
\item Sonstige Skripte, Datenbankdefinitionen oder Templates werden
  unter
\begin{verbatim}
/opt/openbib/
\end{verbatim}
  installiert.
\end{itemize}

\chapter{Download der OpenBib-Software}

Die aktuellste Version der OpenBib-Software k"onnen Sie "uber den
CVS-Server von BerliOS erhalten. Dazu w"ahlen Sie sich ein lokales
Verzeichnis aus, in welches die Dateien abgelegt werden sollen. Zur
Veranschaulichung w"ahlen wir hier das Verzeichnis \texttt{/opt/store}.

Sie wechseln nun in dieses Verzeichnis und f"uhren folgende Befehle
aus: 

\begin{verbatim}
cvs -d:pserver:anonymous@cvs.openbib.berlios.de:/cvsroot/openbib login
\end{verbatim}

Wenn Sie nun nach einen Passwort gefragt werden dr"ucken Sie einfach
die Return-Taste.

\begin{verbatim}
cvs -z3 -d:pserver:anonymous@cvs.openbib.berlios.de:/cvsroot/openbib co openbib
\end{verbatim}

Im Verzeichnis \texttt{/opt/store/openbib} finden Sie nun die gesamte
OpenBib-Distribution. Wenn von nun an einzelne Dateien der
Distribution angesprochen werden, beziehen sie sich jeweils auf dieses
OpenBibRoot-Verzeichnis, also z.B.~\texttt{LICENCE} auf \texttt{/opt/store/openbib/LICENCE}.

Der CVS-Baum ist entsprechend der einzelnen Komponenten in
Verzeichnissen organisiert:

\begin{description}
\item[conv] In diesem Verzeichnis befinden sich die Programme, die bei
  der automatischen Migration in die OpenBib-Katalog-Datenbanken
  ben"otigt werden. Grunds"atzlich sind die Programme wiederum
  entsprechend der Quell-Systeme (sikis, oai etc.) in
  Unterverzeichnissen zusammengefasst. In jedem dieser
  Unterverzeichnisse sind die ben"otigten Programme, um die Daten aus
  dem Quell-System in ein Meta-Format umzuwandeln. Die
  Umwandlungsprogramme, um aus dem Meta-Format die Einladedateien f"ur
  die OpenBib-Katalog-Datenbanken zu erzeugen, befinden sich im
  zus"atzlichen Unterverzeichnisses \texttt{meta}.
\item[db] In diesem Verzeichnis befinden sich die Definitionen f"ur
  die Tabellenstruktur der einzelnen OpenBib-Datenbanktypen bezogen
  auf ein DBMS. Derzeit wird nur \texttt{mysql} unterst"utzt. Die
  Datenbanktypen sind \texttt{pool} f"ur die Katalog-Datenbanken,
  \texttt{session} f"ur die Session-Datenbanken sowie \texttt{user}
  f"ur die Benutzer-Datenbank.
\item[doc] In diesem Verzeichnis befindet sich die Dokumentation zu
  Openbib in Form von LaTeX-Dokumenten.
\item[portal] In diesem Verzeichnis befinden sich alles, was f"ur die
  Web-Anwendung OpenBib ben"otigt wird. Es sind dies
  \begin{itemize}
  \item Konfigurationen f"ur den Apache-Webserver unter
    \texttt{apache}
  \item Statische Elemente wie Bilder, CSS etc. unter \texttt{htdocs} sowie
  \item Perl-Module und Templates unter \texttt{perl}.
  \end{itemize}
\end{description}

\chapter{Installation der notwendigen Infrastruktur}
Das Recherche-Portal, die zugeh"origen Datenbanken sowie die
Migrations-Skripte setzen eine bestimmte Infrastruktur voraus. Diese
besteht aus Betriebssystem, Server-Prozessen, Programmen und
verschiedenen Perl-Modulen. Bevor daher mit der eigentlichen
Installation von OpenBib begonnen werden kann, m"ussen diese
Komponenten installiert werden.

Die Beispielhafte Installation der Komponenten erfolgt am Beispiel
eines Debian Linux Systems in der Version 3.1 (Codename: Sarge).


\section{Installation des Betriebssystems}
OpenBib setzt ein Unix-artiges Betriebssystem, idealerweise Linux,
voraus. Auf die Installation wird hier nicht weiter eingegangen.


\section{Installation des Apache-Webservers}
Aus Performance-Gr"unden ist die Installation von zwei
Apache-Webservern angeraten. Einen leichtgewichtigen mit
mod\_proxy-Support auf Port 80 und einen gro"sen mit mod\_perl-Support
z.B. auf Port 8008. 

Der erste Webserver (Port 80) ist f"ur das Liefern statischer
HTML-Seiten zust"andig sowie als Accelerator-Proxy f"ur den zweiten
Webserver (Port 8008). Die eigentliche Portal-Anwendung l"auft dann
nur noch auf dem zweiten Webserver unter mod\_perl.


\subsection{Installation des ersten Webservers}

Den ersten Webserver installieren Sie mit dem Kommando

\begin{verbatim}
apt-get install apache
\end{verbatim}

Danach wird der Webserver automatisch gestartet und sie finden das
ServerRoot-Verzeichnis unter \texttt{/etc/apache} mit den Konfigurationsdateien
\texttt{httpd.conf} und \texttt{modules.conf} sowie das \texttt{DocumentRoot}-Verzeichnis unter
\texttt{/var/www}.

Deaktivieren Sie in \texttt{/etc/apache/modules.conf} den eventuell
dort aktivierten mod\_perl Support:

\begin{verbatim}
#LoadModule perl\_module /usr/lib/apache/1.3/mod\_perl.so
\end{verbatim}

Aktiviert werden mu"s hingegen das Modul f"ur den Proxy-Support

\begin{verbatim}
LoadModule proxy\_module /usr/lib/apache/1.3/libproxy.so
\end{verbatim}


Nun k"onnen Sie den Accelerator-Proxy in der Konfigurationsdatei
\texttt{/etc/apache/httpd.conf} eintragen.

Tragen Sie dort unter \texttt{<IfModule mod\_proxy.c>} folgendes f"ur die
Weiterleitung auf den zweiten Webserver ein:

\begin{verbatim}
<IfModule mod\_proxy.c>
# Proxy f"ur Dynamischen Content
ProxyPass        /portal/ http://vollqualifizierter.rechnername:8008/portal/
ProxyPassReverse /portal/ http://vollqualifizierter.rechnername:8008/portal/

ProxyReceiveBufferSize 16384

# Einrichtung eines Caches
CacheRoot /usr/db/apache-cache
CacheSize 50000
CacheGcInterval 1
CacheMaxExpire 72
CacheLastModifiedFactor 0.1
CacheDefaultExpire 24

</IfModule>
\end{verbatim}

Als 'vollqualifizierter.rechnername' tragen Sie den DNS-Namen Ihres
Rechners ein, unter dem das Recherche-Portal laufen soll.

Ebenso m"ussen Sie ein CacheRoot-Verzeichnis anlegen, das vom Proxy
verwendet wird und ihm die Ownership des WWW-Benutzers (www-data)
geben. Die geschieht f"ur den o.g. Pfad z.B. mit:

\begin{verbatim}
mkdir -p /usr/db/apache-cache
chown www-data.www-data /usr/db/apache-cache
\end{verbatim}

Die beiden ProxyPass*-Direktiven sorgen daf"ur, da"s alle Anfragen,
die unter dem Basis-URL \texttt{/portal} liegen an den zweiten Webserver
weitergeleitet werden.

Nach einem Restart des ersten Webservers ist dieser eingerichtet.

\begin{verbatim}
/etc/init.d/apache restart
\end{verbatim}


Hinweis: Falls Sie VirtualHosts verwenden, so k"onnen Sie dort auch
den Abschnitt mit \texttt{<IfModule mod\_proxy.c>} f"ur die jeweiligen Hosts
unterbringen.


\subsection{Installation des zweiten Webservers}

Den zweiten Webserver installieren Sie mit dem Kommando

\begin{verbatim}
apt-get install apache-perl
\end{verbatim}

Danach wird der Webserver automatisch gestartet und sie finden das
Server-Root-Verzeichnis unter \texttt{/etc/apache-perl} mit der
Konfigurationsdatei \texttt{httpd.conf} sowie das
\texttt{DocumentRoot}-Verzeichnis ebenfalls unter \texttt{/var/www}.

"Andern Sie nun in \texttt{/etc/apache-perl/httpd.conf} den Port auf
8008:

\begin{verbatim}
Port 8008
\end{verbatim}

Tragen Sie nun folgende Direktive in \texttt{/etc/apache-perl/httpd.conf} ein:

\begin{verbatim}
Include /etc/apache-perl/openbib.conf
\end{verbatim}

Gegebenenfalls mu"s diese Direktive im entsprechenden
\texttt{VirtualHosts}-Block eingetragen werden.

Mit dem Include werden die Konfigurationen f"ur das Recherche-Portal
eingeladen. Kopieren sie daf"ur nun die Datei
\texttt{portal/apache/openbib.conf} aus der OpenBib-Distribution in
das Verzeichnis \texttt{/etc/apache-perl}.

\begin{verbatim}
cp /opt/store/openbib/portal/apache/openbib.conf /etc/apache-perl/
\end{verbatim}


\section{Installation und Konfiguration des MySQL-DBMS}
\label{sec:instmysql}


\subsection{Installation}

OpenBib verwendet verschiedene Datenbanken. So gibt es auf jedem
Rechner eine session-Datenbank, in die Sitzung-Informationen sowie
generelle Informationen zur Konfiguration gespeichert werden. Auf
genau einem Rechner (wenn man die Lastverteilung verwendet) existiert
eine user-Datenbank, in der Benutzer-Daten und Einstellungen
abgespeichert werden. Und schlie"slich gibt es auf jedem Rechner
mindestens eine, normalerweise aber unz"ahlige, pool-Datenbanken, in
denen die eigentlichen Kataloge untergebracht werden in der Form: eine
Datenbank f"ur einen Katalog.

Alle diese Datenbanken werden mit dem DBMS MySQL verwaltet. MySQL
installieren Sie mit:

\begin{verbatim}
apt-get install mysql-server
\end{verbatim}

Mit diesem Paket werden verschiedene andere Pakete, die noch ben"otigt
werden, automatisch mitinstalliert, z.B. die Perl-Module, um auf die
MySQL-Datenbanken zugreifen zu k"onnen.

Nach der Installation wird der MySQL-Server automatisch gestartet.

Setzen Sie nun ein Passwort \texttt{<Passwort>} f"ur den Benutzer root mit

\begin{verbatim}
mysqladmin -u root -h localhost password <Passwort>
\end{verbatim}

Falls Sie sich 'ausgesperrt' haben sollten, so k"onnen Sie den
\texttt{mysqld} mit der Option
\begin{verbatim}
--skip-grant-tables
\end{verbatim}
neu starten, den Eintrag f"ur root l"oschen, \texttt{mysqld} restarten
und dann das Passwort wie beschrieben neu setzen.

F"ur die verschiedenen Datenbanken von OpenBib ist es am einfachsten,
diesen Benutzer \texttt{root} und das zugeh"orige Passwort
\texttt{<Passwort>} f"ur den Zugriff zu verwenden. Aus
Sicherheitsgr"unden -- insbesondere wenn weitere
Produktionsdatenbanken auf dem lokalen Rechner unter MySQL betrieben
werden -- kann es jedoch empfehlenswert sein, einen weniger
privilegierten Benutzer, z.B. \texttt{openbib}, hierf"ur zu verwenden.
Als Beispielpasswort verwenden wir im Folgenden \texttt{TopSecret}.


\begin{verbatim}
openbib:/opt/store/openbib/tools# mysql -u root -p mysql
Enter password:
Reading table information for completion of table and column names
You can turn off this feature to get a quicker startup with -A

Welcome to the MySQL monitor.  Commands end with ; or \g.
Your MySQL connection id is 2857 to server version: 4.0.22_Debian-6-log

Type 'help;' or '\h' for help. Type '\c' to clear the buffer.


mysql> grant usage on *.* to openbib@localhost identified by 'TopSecret';
Query OK, 0 rows affected (0.04 sec)

mysql> flush privileges;
Query OK, 0 rows affected (0.04 sec)
\end{verbatim}

Mit diesem \texttt{grant}-Befehl in Mysql wird in der
\texttt{user}-Tabelle der Datenbank \texttt{mysql} der Benutzer
\texttt{openbib} f"ur \texttt{localhost} angelegt, jedoch zun"achst
ohne weitere Rechte. Diese m"ussen nun h"andisch gesetzt werden:

\begin{verbatim}
mysql> update user set Select_priv='Y', Insert_priv='Y',
 Update_priv='Y', Delete_priv='Y', Create_priv='Y', Drop_priv='Y', 
 Index_priv='Y', Alter_priv='Y', File_priv='Y', Reload_priv='Y' 
 Lock_tables_priv='Y', Create_tmp_table_priv='Y' where User='openbib';
mysql> flush privileges;
\end{verbatim}

Die Privilegien \texttt{Shutdown\_priv},
\texttt{Process\_priv} sowie \texttt{Grant\_priv}
und \texttt{References\_priv} werden nicht ben"otigt uns sollten auch
aus Sicherheitsgr"unden nicht vergeben werden.

Prinzipiell ist es in OpenBib m"oglich f"ur alle drei
Datenbank-Klassen (pool, session und user) einen verschiedenen
Benutzerzugang zu konfigurieren. Der Einfachheit halber wird im
Folgenden jeweils nur einer, konkret der gerade eingerichtete Benutzer
\texttt{openbib} verwendet.


\subsection{Konfiguration}

Die Konfiguration bezieht sich im Wesentlichen auf die Datei
\texttt{/etc/mysql/my.cnf}. Dort m"ussen verschiedene Einstellungen
vorgenommen werden.

F"ur die Volltextsuche sind dies:

\begin{description}
\item[ft\_min\_word\_len] In den Sektionen \texttt{[mysqld]} und
  \texttt{[myisamchk]} mu"s die minimale L"ange der zu indizierenden
  Worte angegeben werden. Der Standardwert bei Weglassung der
  Einstellung ist 3. Oft ist es jedoch sinnvoll diesen Wert auf 2 oder
  1 herunterzusetzen, damit weitere Kurzw"orter -- dann jedoch auf
  Kosten eines wachsenden Index -- recherchierbar werden.
\begin{verbatim}
# Minimale Wortlaenge fuer Volltextsuche
ft_min_word_len = 2
\end{verbatim}
\item[ft\_stopword\_file] In der Sektion \texttt{[mysqld]} sollte
  durch Angabe dieses Parameters \textbf{ohne Wert} die
  Ber"ucksichtigung von Stopwortlisten unterdr"uckt werden.
\begin{verbatim}
# Stoppworte deaktivieren
ft_stopword_file =
\end{verbatim}
\end{description}

\chapter{Installation des OpenBib Web-Portals}

Die Installation von OpenBib besteht aus

\begin{itemize}
\item Installation und Konfiguration der Perl-Module des Web-Portals
\item Installation der Konvertierungs-Infrastruktur
\item Anlegen und Konfiguration der Datenbanken
\end{itemize}

\section{Installation der Perl-Module des Web-Portals}

Erzeugen Sie zun"achst das Verzeichnis \texttt{/usr/local/lib/site\_perl}, das
zum Standard-Include-Pfad von Perl geh"ort, falls dies noch nicht
existiert:

\begin{verbatim}
mkdir /usr/local/lib/site_perl
\end{verbatim}

Kopieren Sie nun aus dem CVS-Baum unter \texttt{/opt/store/openbib}
das Verzeichnis 

\begin{verbatim}
portal/perl/modules/OpenBib
\end{verbatim}

in das gerade angelegte Verzeichnis:

\begin{verbatim}
cp -a /opt/store/openbib/portal/perl/modules/OpenBib /usr/local/lib/site_perl/
\end{verbatim}

Diese Perl-Module ben"otigen weitere Perl-Module, die als
Debian-Pakete nachinstalliert werden m"ussen:

\begin{itemize}
\item Das Perl Template Toolkit inkl. Dokumentation
\begin{verbatim}
apt-get install libtemplate-perl libtemplate-perl-doc
\end{verbatim}
  Zus"atzlich wird das hierzu ben"otigte Paket libappconfig-perl mit
  installiert.
\item Apache::Request
\begin{verbatim}
apt-get install libapache-request-perl
\end{verbatim}
\item Apache::Reload
\begin{verbatim}
apt-get install libapache-reload-perl
\end{verbatim}
\item LWP
\begin{verbatim}
apt-get install libwww-perl
\end{verbatim}
\item DBI und DBD::Mysql
\begin{verbatim}
apt-get install libdbi-perl libdbd-mysql-perl
\end{verbatim}
\item Log4Perl
\begin{verbatim}
apt-get install liblog-log4perl-perl libipc-shareable-perl
\end{verbatim}
\item MIME::Lite
\begin{verbatim}
apt-get install libmime-lite-perl
\end{verbatim}
\item Email::Valid
\begin{verbatim}
apt-get install libemail-valid-perl
\end{verbatim}
\item oai-perl 2.15
Dieses Modul muss h"andisch von
\texttt{http://oai-perl.sourceforge.net/} installiert werden.
\end{itemize}

\section{Konfiguration des Web-Portals}
\label{sec:confportal}

Die zentrale Konfigurationsdatei in OpenBib ist das Modul
OpenBib::Config. Die zugeh"orige Datei befindet sich in Form einer
anzupassenden Beispielkonfiguration unter 

\begin{verbatim}
/usr/local/lib/site_perl/OpenBib/Config.pm-dist
\end{verbatim}

Sie mu"s nun nach Config.pm kopiert 

\begin{verbatim}
cd /usr/local/lib/site_perl/OpenBib/
cp -a Config.pm-dist Config.pm
\end{verbatim}

und angepa"st werden.


\subsection{Server-Einstellungen}

\begin{description}
\item[servername] Geben Sie hier den Namen des lokalen Rechners
  ein. 
\begin{verbatim}
servername => 'openbib.foo.bar',
\end{verbatim}

\item[loadbalancerservername] Geben Sie hier den Namen des Rechners
  ein, der die Lastverteilung "ubernimmt. Falls insgesamt nur ein
  Rechner zur Verf"ugung steht, dann entspricht der Eintrag dem von
  \texttt{servername}.
\begin{verbatim}
loadbalancerservername => 'openbib.foo.bar',
\end{verbatim}

\item[loadbalancertargets] Geben Sie hier die Namen der Rechner ein,
  auf die "uber die die Last verteilt werden soll. Falls insgesamt nur
  ein Rechner zur Verf"ugung steht, dann entspricht der Eintrag dem
  von \texttt{servername}.
\begin{verbatim}
loadbalancertargets => [
                        'openbib.foo.bar',
                       ],
\end{verbatim}
\end{description}


\subsection{Datenbank-Einstellungen}


\begin{description}
\item[dbms] An dieser Stelle wird das verwendete DBMS
  konfiguriert. Derzeit kann dort nur \texttt{mysql} verwendet
  werden. Daher mu"s dieser Eintrag nicht ge"andert werden.
\begin{verbatim}
dbms      => 'mysql',
\end{verbatim}

\item[dbimodule] An dieser Stelle wird der Name des verwendeten
  DBI-Modules konfiguriert. Derzeit kann dort nur \texttt{mysql}
  verwendet werden. Daher mu"s dieser Eintrag nicht ge"andert werden.
\begin{verbatim}
dbimodule => 'mysql',
\end{verbatim}
\item[dbhost] An dieser Stelle wird der Name des Rechners
  konfiguriert, auf dem die Katalog-Datenbanken liegen. Dies sollte
  immer der lokale Rechner sein. Daher mu"s dieser Eintrag nicht
  ge"andert werden.
\begin{verbatim}
dbhost    => 'localhost',
\end{verbatim}
\item[dbport] An dieser Stelle wird der Port festgelegt, "uber den auf
  die Katalog-Datenbanken auf dem Server \texttt{dbhost} zugegriffen
  werden kann. Mysql verwendet standardm"a"sig \texttt{3306}. Daher
  mu"s dieser Eintrag nicht ge"andert werden.
\begin{verbatim}
dbport    => '3306',
\end{verbatim}
\item[dbuser] "Uber diesen Benutzer wird auf die Katalog-Datenbanken
  zugegriffen. Verwenden Sie aus Sicherheitsgr"unden einen Benutzer,
  der nicht identisch mit dem \texttt{root}-Benutzer ist, da dieser
  die volle Kontrolle "uber das Mysql-DBMS besitzt (s.o.). Verwenden
  Sie hier entsprechend \ref{sec:instmysql} den Benutzer
  \texttt{openbib}. Daher mu"s dieser Eintrag entsprechend ge"andert
  werden.
\begin{verbatim}
dbuser    => 'openbib',
\end{verbatim}
\item[dbpasswd] Tragen Sie hier das Passwort f"ur den \texttt{dbuser}
  ein. Entsprechend \ref{sec:instmysql} ist dies beispielhaft
  \texttt{TopSecret}. Daher mu"s dieser Eintrag entsprechend ge"andert
  werden.
\begin{verbatim}
dbpasswd  => 'TopSecret',
\end{verbatim}
\end{description}

Entsprechend \texttt{dbuser} und \texttt{dbpasswd} werden nun die
Parameter \texttt{sessiondbuser}, \texttt{sessiondbpasswd} sowie
\texttt{userdbuser} und \texttt{userdbpasswd} gesetzt. Bei den
Parametern f"ur \texttt{*host}, \texttt{*port} und \texttt{*name}
k"onnen die Standardwerte "ubernommen werden.



\subsection{Administrations-Einstellungen}
\label{ssec:admineinstellung}

\begin{description}
\item[adminuser] Mit diesem Benutzernamen greifen Sie sp"ater auf das
  Web-Administrations-Interface zu. "Andern Sie diesen z.B. auf
  \texttt{admin}.
\begin{verbatim}
adminuser     => 'admin',
\end{verbatim}
\item[adminpasswd] Mit diesem Passwort authentifizieren Sie sich
  sp"ater als \texttt{adminuser}.
\begin{verbatim}
adminpasswd   => 'NochGeheimer',
\end{verbatim}
\end{description}


\subsection{Organisationsstruktur der Kataloge}

\begin{description}
\item[units] Mit dem Parameter \texttt{units} legen Sie die
  Organisationsstruktur Ihrer Kataloge fest. Traditionell werden an
  einer Universit"at z.B. die Kataloge nach Fakult"aten
  organisiert. F"ur jede Organisationseinheit definieren Sie hier
  einen Kurz-Eintrag \texttt{short}, in den z.B. durch vorgestellte
  Zahlen eine Reihenfolge gebracht werden kann, und eine Beschreibung
  \texttt{desc}. 
\begin{verbatim}
units => [
          {
           "short","0ungeb",
           "desc","Fakult&auml;tsungebunden",
          },
          {
           "short","1wiso",
           "desc","Wirtschafts- u. Sozialwissenschaftliche Fakult&auml;t",
          },
          {
           "short","2recht",
           "desc","Rechtswissenschaftliche Fakult&auml;t",
          },
          {
           "short","3ezwheil",
           "desc","Erziehungswiss. u. Heilp&auml;d. Fakult&auml;t",
          },
          {
           "short","4phil",
           "desc","Philosophische Fakult&auml;t",
          },
          {
           "short","5matnat",
           "desc","Mathematisch-Naturwissenschaftliche Fakult&auml;t"
,
          },
         ],
\end{verbatim}
\end{description}

Dieser Parameter wird leider noch nicht in allen Modulen verwendet, so
da"s Sie diese h"andisch entsprechend "andern m"ussen.
 
\subsection{URI-Konfiguration}

In OpenBib k"onnen Sie f"ur jede Web-Komponente festlegen, unter
welchem URI sie angesprochen werden soll. Die OpenBib-interne
Verlinkung ist damit entsprechend konsistent. Alle zugeh"origen
Parameter enden auf \texttt{\_loc}. Im Normalfall brauchen Sie diese
nicht zu "andern. Standardm"a"sig beginnen die URI's mit \texttt{/portal/}.



\subsection{Verschiedene Konfigurationsparameter}

\begin{description}
\item[mail\_prog] Hier geben Sie den Pfad zu dem Program an, mit dem
  das System Mails verschickt.
\item[admin\_email] Hier geben Sie die E-Mail-Adresse an, die bei Mails
  an den Administrator, also Sie, verwendet werden soll.
\item[contact\_email] Hier geben Sie die E-Mail-Adressen an, die als
  generelle Kontakt-Adresse auf den Web-Seiten genannt werden soll.
\item[hbzmonofernleih\_exturl] Hier geben Sie den Basis-URL an, der bei
  einem Sprung in die Mono\-graphien-Fernleihe der Digitalen Bibliothek
  NRW verwendet werden soll.
\item[hbzzeitfernleih\_exturl] Hier geben Sie den Basis-URL an, der bei
  einem Sprung in die Aufsatz-Fernleihe der Digitalen Bibliothek
  NRW verwendet werden soll.
\item[ezb\_exturl] Hier geben Sie den Basis-URL an, der bei
  einem Sprung in die Elektronische Zeitschriftenbibliothek (EZB) in
  Regensburg verwendet werden soll.
\item[dbis\_exturl] Hier geben Sie den Basis-URL an, der bei einem
  Sprung in das DatenBankInformationsSystem (DBIS)  in Regensburg
  verwendet werden soll. 
\item[databasechoice\_maxcolumn] Hier legen Sie fest, auf wieviele
  Spalten die Kataloge in der Datenbankauswahlseite verteilt werden
  sollen. 
\item[web\_seed] Hier legen sie ein 'geheimes' Wort fest, welches
  zusammen mit dem Katalog-Datenbanknamen per MD5 ein
  Web-Zugriffspasswort ergibt, mit dem automatisiert auf entfernt
  exportierte Katalogdaten per HTTP zugegriffen werden kann. Dieser
  Parameter wird bald obsolet, da die Zugriffsinformationen in naher
  Zukunft direkt aus der session-Datenbank geholt werden.
\item[benchmark] Hier legen Sie fest, ob Benchmark-Informationen
  ausgegeben werden sollen. Derzeit f"uhrt dies zu Fehlern und sollte
  daher deaktiviert werden. Das ist auch die Standardeinstellung.
\item[testsystem] Hier legen Sie fest, ob es sich bei dem Rechner um
  ein Testsystem handelt. Es wird dann zus"atzlich eine entsprechende
  Informationsmeldung in die Basis-Webseite eingeblendet.
\end{description}


\subsection{Template-Konfiguration}

In OpenBib werden bei fast allen Modulen Templates zur Trennung von
Programm- und HTML-Code verwendet. An dieser Stelle ist von Ihnen nur
der Pfad zu dem Verzeichnis zu setzen, in dem Sie die Templates
ablegen.

\begin{description}
\item[tt\_include\_path] Setzen Sie diesen Pfad, z.B. auf
  \texttt{/opt/openbib/templates}.
\end{description}

Alle zugeh"origen Parameter zur Namensdeklaration einzelner Templates
enden auf \texttt{\_ttname}. Im Normalfall brauchen Sie diese nicht
"andern.


\subsection{Generelle Pfad-Konfiguration}

Mit den auf \texttt{\_dir} endenden Parameter wird die generelle
Pfad-Struktur definiert. Bei einer Installation unter
\texttt{/opt/openbib} brauche Sie hier nichts "andern.

\begin{description}
\item[base\_dir] Dies ist das Basis-Verzeichnis von OpenBib.
\item[conv\_dir] In diesem Verzeichnis befinden sich die Programme zur
  Konvertierung.
\item[autoconv\_dir] In diesem Verzeichnis befindet sich die
  Infrastruktur, um automatisiert Kataloge in OpenBib
  einzuspielen. Dazu geh"ort ein Unterverzeichnis \texttt{bin} f"ur
  die Steuer-Programme, ein Unterverzeichnis \texttt{pools} f"ur eine
  lokale Kopie der jeweiligen Katalogdaten und ein Unterverzeichnis
  \texttt{data} f"ur tempor"are Dateien, die beim Einspielen erzeugt
  werden.
\item[dbdesc\_dir] In diesem Verzeichnis befinden sich die
  Datenbankdefinitionen nach Datenbanksystem getrennt. Derzeit
  existieren nur Definitionen f"ur MySQL.
\item[tool\_dir] In diesem Verzeichnis befinden sich Hilfs-Skripte,
  z.B. zum Anlegen einer Datenbank.
\end{description}


\section{Anlegen der Verzeichnis-Struktur}

Als n"achstes wird die in \ref{sec:confportal} definierte Verzeichnis
struktur erzeugt.

\begin{verbatim}
mkdir /opt/openbib
mkdir /opt/openbib/bin
mkdir /opt/openbib/autoconv
mkdir /opt/openbib/autoconv/bin
mkdir /opt/openbib/autoconv/data
mkdir /opt/openbib/autoconv/pools
mkdir /opt/openbib/conv
mkdir /opt/openbib/db
mkdir /opt/openbib/templates
\end{verbatim}

Alle diese Verzeichnisse m"ussen f"ur den Prozess des Web-Servers
les- und schreibbar sein. Daher m"ussen Sie Ihre Besitzer entsprechend
"andern.

\begin{verbatim}
chown -R www-data.www-data /opt/openbib
\end{verbatim}

\section{Installation der Templates des Web-Portals}

Um Programm-Code und HTML-Code des Web-Portals zu trennen, wurde in
vielen OpenBib-mod\_perl-Modulen das Perl Template
Toolkit\footnote{\texttt{http://www.template-toolkit.org/}}
eingesetzt.

Die Templates werden entsprechend der Konfiguration in
\ref{sec:confportal} unter \texttt{/opt/openbib/templates} abgelegt.

\begin{verbatim}
cp /opt/store/openbib/portal/perl/templates/* /opt/openbib/templates/
\end{verbatim}

\section{Installation der Datenbankdefinitionen}

Die Datenbankdefinitionen werden entsprechend der Konfiguration in
\ref{sec:confportal} unter \texttt{/opt/openbib/db} abgelegt.

\begin{verbatim}
cp -r /opt/store/openbib/db/mysql /opt/openbib/db/
\end{verbatim}

\section{Installation verschiedener Hilfs-Skripte}

F"ur verschiedene Aufgaben stehen Hilfs-Skript bereit, die
entsprechend der Konfiguration von \texttt{tool\_dir} in
\ref{sec:confportal} unter \texttt{/opt/openbib/bin} abgelegt werden.
So z.B. zum Anlegen der Benutzerdatenbank, Sessiondatenbank usw. Im
CVS-Baum finden Sie diese Skripte unter \texttt{tools}.

\begin{verbatim}
cp /opt/store/openbib/tools/* /opt/openbib/bin/
\end{verbatim}

\section{Vorbereiten des Web-Bereichs}

Im \texttt{DocumentRoot}-Verzeichnis muss die notwendige Infrastruktur
installiert werden. Es sind dies

\begin{itemize}
\item die Stylesheets
\item die Bilder und
\item die Hilfeseite
\end{itemize}

Beispiele finden Sie im CVS-Baum unter
\texttt{portal/htdocs/}. Kopieren Sie diese Verzeichnisse in Ihr
\texttt{DocumentRoot}-Verzeichnis.

\begin{verbatim}
cp -r /opt/store/openbib/portal/htdocs/* /var/www/
\end{verbatim}


\section{Installation der Konvertierungsprogramme}

Die ben"otigten Konvertierungsprogramme m"ussen nun aus dem CVS-Baum
unter \texttt{conv} in die zugeh"origen Verzeichnisse unter
\texttt{/opt/openbib} kopiert werden.

\begin{verbatim}
cp /opt/store/openbib/conv/meta/* /opt/openbib/conv/
cp /opt/store/openbib/conv/sikis/autoconv-sikis.pl /opt/openbib/autoconv/bin/
\end{verbatim}

\section{Anlegen der Benutzer- und Session-Datenbanken}

Grundlage f"ur den Betrieb des Web-Portals ist die Existenz der
Benutzer- und Session-Datenbanken. Diese werden nun mit den
Hils-Skripten \texttt{createuser.pl} und \texttt{createsession.pl} aus
\texttt{/opt/openbib/bin} erzeugt.

\begin{verbatim}
openbib:/opt/openbib/bin# ./createuser.pl
openbib:/opt/openbib/bin# ./createsession.pl
\end{verbatim}

Re-Starten Sie nun die Webserver. Es werden dann entsprechend der
Konfiguration in \texttt{openbib.conf} die Module des Portals in den
Webserver eingeladen.

Damit sind alle Vorarbeiten f"ur den Betrieb des Portal vollzogen --
bis auf die Existenz der eigentlichen Kataloge -- und "uber den
Web-Browser kann nun bereits darauf zugegriffen werden. Die
erfolgreiche Installation des Portals dokumentiert Abbildung
\ref{bild:fertig} mit derzeit 0 Katalog-Datenbanken.

\begin{figure}
\begin{shadowenv}
  \vspace{4mm}
    \centering \begin{minipage}[b]{1.0\textwidth}
      \centering \includegraphics[width=14cm]{fertig1.eps}
    \end{minipage}
    \caption{Nach Anlegen der Benutzer- und Session-Datenbank und
      Restart des Web-Servers erscheint das Portal in seiner Standardkonfiguration}
  \label{bild:fertig}
  \vspace{3mm}
\end{shadowenv}
\end{figure}


\chapter{Installation von Katalogen und automatischer Konvertierung}

Nachdem nun das Web-Portal installiert und konfiguriert ist, k"onnen
Kataloge definiert werden.

Generell wird davon ausgegangen, dass die Katalogdaten aus externen
(Bibliotheks-)Systemen stammen, wie z.B. aus SunRise der Firma
Sisis\footnote{Sisis und Sunrise sind eingetragene Marken der Firma
  Sisis Informationssysteme GmbH, M"unchen.} oder OAI-Datenquellen.

Daher gestaltet sich die Installation eines Kataloges typischerweise
aus mehreren Schritten:

\begin{enumerate}
\item Der Katalog wird mit Meta-Daten "uber das
  Administrations-Interface des Web-Portals erfasst und
  angelegt. Ebenso wird die Quelle der (exportierten) Daten definiert.
\item Auf dem Rechner des Bibliothekssystem wird gegebenenfalls eine
  regelm"a"sige Entlade- und/oder Konvertierungsroutine installiert,
  die die Bibliotheksdaten z.B. in einem (mit Passwort gesch"utzten)
  Web-Bereich ablegt.
\end{enumerate}

Welche der Schritte in welchem Umfang vollzogen werden, h"angt
ma"sgeblich von der Art des Kataloges ab.

Im Folgenden werden diese Schritte anhand eines Beispielkataloges
durchgef"uhrt, dessen exportierte Daten bereits auf der
OpenBib-Website vorliegen. Auf diese Weise kann anhand dieses
Beispielkataloges bereits das Web-Portal evaluiert werden.

\section{Aufnahme eines Katalogs "uber das Administrations-Interface}

Um einen neuen Katalog mit seinen Meta-Daten zu erfassen wird das
Administrations-Interface des Web-Portals genutzt.

Rufen Sie daher den folgenden URL entsprechend der bereits erfolgten
Installation und Konfiguration auf:

\begin{verbatim}
http://openbib.foo.bar/portal/admin
\end{verbatim}

In der Eingabemaske authentifizieren Sie sich mit dem unter
\ref{ssec:admineinstellung} definierten \texttt{adminuser} 'admin' und
dem \texttt{adminpasswd} 'NochGeheimer'.

Nach erfolgreicher Authentifikation erscheint die Erfolgsmeldung
\emph{Sie sind nun als Administrator eingeloggt.} sowie ein Menu mit
den Eintr"agen \texttt{Kataloge}, \texttt{Views},
\texttt{Sessions} und \texttt{Import}.

\begin{description}
\item[Kataloge] Hier werden neue Kataloge mit Ihren Meta-Daten
  aufgenommen, ge"andert und gel"oscht. Speziell werden hier auch die
  Quellen der Katalogdaten konfiguriert.
\item[Views] Hier k"onnen Katalog/Instituts-Sichten (Views) definiert
  werden. Auf diese Weise k"onnen aus allen Katalogen des
  Recherche-Portals Untermengen, z.B. thematisch, mit eigenem Namen
  und Beschreibung definiert werden, so da"s direkt in diese
  Vorauswahl an Katalogen in das Portal hereingesprungen werden kann.
\item[Sessions] Hier finden Sie eine Aufstellung der existierenden
  Sessions. Zu jeder Session k"onnen Sie sich weitergehende
  Informationen anzeigen lassen, z.B. welche Recherchen mit welchen
  Suchbegriffen in der betreffenden Session get"atigt wurden.
\item[Import] Hier k"onnen Sie automatische Importe f"ur einzelne
  Kataloge veranlassen.
\end{description}

\end{document}
